\documentclass[preprint,12pt]{elsarticle}
\usepackage{amsmath,amssymb,amsfonts,amsthm}
\usepackage{algorithm}
\usepackage{algpseudocode}
\usepackage{graphicx}
\usepackage{xcolor}
\usepackage{booktabs}
\usepackage{adjustbox}
\emergencystretch=3em
\usepackage{tikz}
\usetikzlibrary{arrows.meta,backgrounds,fit}
\usepackage{pgfplots}
\pgfplotsset{compat=1.18}
\usepgfplotslibrary{groupplots}
\usepackage{xurl}
\usepackage[hidelinks]{hyperref}

\newtheorem{proposition}{Proposition}
% generated by paper/tools/make_paper_assets.py
\newcommand{\nobBFourchg}{11.7}
\newcommand{\nobLiRAchg}{0.27}
\newcommand{\nobY}{97.7}
\newcommand{\nobZ}{97.7}
\newcommand{\nobX}{100.0}
\newcommand{\nobhlBFourchg}{1.76}
\newcommand{\nobhlLiRAchg}{0.80}
\newcommand{\nobhlY}{52}
\newcommand{\nobhlZ}{75}
\newcommand{\abiBFourchg}{0.019}
\newcommand{\abiLiRAchg}{0.000}
\newcommand{\abiSav}{24.0}
\newcommand{\geaBFourchg}{0.000}
\newcommand{\geaLiRAchg}{0.000}
\newcommand{\geaSav}{29.8}
\newcommand{\nobSav}{27.6}
\newcommand{\nobSnapChg}{63.7}
\newcommand{\nobBsixChg}{0.13}
\newcommand{\nobhlJBFour}{890.7}
\newcommand{\nobhlJBSix}{887.0}
\newcommand{\nobhlJBSeven}{890.7}
\newcommand{\nobhlJLiRA}{880.9}
\newcommand{\nobhlJOstar}{880.9}
\newcommand{\nobhlBsixChg}{1.07}
\newcommand{\nobhlPbsix}{0.31}
\newcommand{\nobPbfour}{0.023}
\newcommand{\nobPbsix}{0.88}
\newcommand{\globalTests}{36}
\newcommand{\globalSig}{0}
\newcommand{\globalMinP}{0.28}
\newcommand{\nobstPbsix}{0.64}
\newcommand{\nobstBsixDays}{25}
\newcommand{\dpMilpDiff}{1\times10^{-14}}
\newcommand{\abiTestDays}{41.8}
\newcommand{\geaTestDays}{28.1}
\newcommand{\nobTestDays}{7.5}
\newcommand{\ilpGermany}{1.6}
\newcommand{\ilpTA}{11.4}
\newcommand{\ilpGeant}{4.1}
\newcommand{\safetyTA}{0.36}
\newcommand{\nobIlpMs}{537}
\newcommand{\geaIlpS}{3.5}
\newcommand{\nobGaS}{28}
\newcommand{\geaGaS}{42}
\newcommand{\abiLatBFour}{0.99}
\newcommand{\abiLatLiRA}{0.008}
\newcommand{\geaLatBFour}{1.60}
\newcommand{\geaLatLiRA}{1.571}
\newcommand{\nobLatBFour}{0.95}
\newcommand{\nobLatLiRA}{0.008}
\newcommand{\abiIntervals}{48\,108}
\newcommand{\geaIntervals}{10\,784}
\newcommand{\nobIntervals}{8\,640}
\newcommand{\emuTrans}{18}
\newcommand{\emuRuleMs}{99}
\newcommand{\emuLossWin}{0.22}
\newcommand{\emuLossAll}{0.014}
\newcommand{\emuTransS}{1.11}
\newcommand{\emuLiRALoss}{0.000}
\newcommand{\emuLiRAtrans}{0}
\newcommand{\emuLostPerTrans}{3.7}
\newcommand{\emuLinkSave}{31}
\newcommand{\emuProbePkts}{460\,800}
\newcommand{\emuHaveData}{1}
\newcommand{\modXgbChg}{12.8}
\newcommand{\modXgbJ}{937.2}
\newcommand{\modLgbmChg}{10.1}
\newcommand{\modLgbmJ}{924.1}
\newcommand{\modRfChg}{0.5}
\newcommand{\modRfJ}{877.2}
\newcommand{\modMlpChg}{44.8}
\newcommand{\modMlpJ}{1093.7}
\newcommand{\ablStressO}{894.3}
\newcommand{\ablStressFull}{897.4}
\newcommand{\ablStressNoPrev}{902.0}
\newcommand{\ablStressNoCS}{898.6}
\newcommand{\ablStressNoDagger}{899.0}
\newcommand{\ablStressNoLag}{895.6}
\newcommand{\ablStressNoSafe}{889.9}
\newcommand{\ablStressNoSafeViol}{0.17}
\newcommand{\ablHlO}{880.9}
\newcommand{\ablHlFull}{880.9}
\newcommand{\ablHlNoPrev}{880.9}
\newcommand{\ablHlNoCS}{880.9}
\newcommand{\ablHlNoDagger}{880.9}
\newcommand{\ablHlNoLag}{880.9}
\newcommand{\ablHlNoSafe}{877.3}
\newcommand{\ablHlNoSafeViol}{0.14}
\newcommand{\KBFouratEight}{939.3}
\newcommand{\KBsixAtEight}{875.3}
\newcommand{\KBFouratSixFour}{935.7}
\newcommand{\KBsixAtSixFour}{875.3}
\newcommand{\KactualMax}{51}
\newcommand{\thSixBFourchg}{10.4}
\newcommand{\thSixLiRAchg}{0.40}
\newcommand{\driftFaAbi}{0.17}
\newcommand{\driftThrAbi}{200}
\newcommand{\driftFaGea}{0.00}
\newcommand{\driftThrGea}{200}
\newcommand{\driftFaNob}{0.93}
\newcommand{\driftThrNob}{10}
\newcommand{\foldNobBFour}{371.5}
\newcommand{\foldNobBSix}{350.1}
\newcommand{\foldNobBSeven}{357.9}
\newcommand{\foldNobLiRA}{350.3}
\newcommand{\foldNobO}{350.0}
\newcommand{\nobAFfour}{0.855}
\newcommand{\nobAFlira}{0.846}
\newcommand{\nobAFsnap}{0.895}
\newcommand{\nobBFourtrans}{23.4}
\newcommand{\nobLiRAtrans}{0.53}
\newcommand{\nobJfour}{931.7}
\newcommand{\nobJsix}{875.3}
\newcommand{\nobJlira}{875.9}
\newcommand{\nobJstar}{874.7}
\newcommand{\nobJzero}{2163.7}
\newcommand{\wearBigChiBFour}{0.94}
\newcommand{\wearBigChiLiRA}{0.85}
\newcommand{\wearCrossBFour}{none}
\newcommand{\liraConstRuns}{30}
\newcommand{\liraTotalRuns}{60}
\newcommand{\gainChanges}{96}
\newcommand{\gainLowShare}{100}
\newcommand{\gainMedian}{0.00}
\newcommand{\gainTenthShare}{99}
\newcommand{\gainMax}{0.41}
\newcommand{\gainMedianBelow}{0.01}
\newcommand{\caseBFourchg}{88}
\newcommand{\caseLiRAchg}{2}
\newcommand{\boundabibaseGap}{0.000}
\newcommand{\boundabibaseLoss}{0.000}
\newcommand{\boundgeabaseGap}{0.011}
\newcommand{\boundgeabaseLoss}{0.011}
\newcommand{\boundnobbaseGap}{0.024}
\newcommand{\boundnobbaseLoss}{0.001}
\newcommand{\boundnobhlGap}{0.463}
\newcommand{\boundnobhlLoss}{0.043}
\newcommand{\boundnobstGap}{0.554}
\newcommand{\boundnobstLoss}{0.045}
\newcommand{\boundMaxLoss}{0.045}
\newcommand{\boundMaxGap}{0.55}
\newcommand{\boundIntervals}{21\,203}
\newcommand{\boundAllOptimal}{all}
\newcommand{\boundMedianMs}{214}
\newcommand{\boundMaxMipGap}{9\times10^{-7}}
\newcommand{\staleRealMaxDJ}{0.000}
\newcommand{\staleRealDJtext}{less than 0.001}
\newcommand{\staleRealMaxViol}{0.008}
\newcommand{\staleNobMaxDJ}{0.002}
\newcommand{\staleNobViolMin}{0.00}
\newcommand{\staleNobViolMax}{0.19}
\newcommand{\staleBzerobase}{0.05}
\newcommand{\staleLiRAViolbase}{0.05}
\newcommand{\staleBzerohl}{0.19}
\newcommand{\staleLiRAViolhl}{0.14}
\newcommand{\staleBzerost}{0.19}
\newcommand{\staleLiRAViolst}{0.17}
\newcommand{\marginSameViol}{yes}
\newcommand{\marginLiRAhlJ}{886.2}
\newcommand{\staleSeeds}{3}
\newcommand{\staticJbase}{875.9}
\newcommand{\staticJhd}{883.7}
\newcommand{\staticJhl}{880.9}
\newcommand{\staticJst}{902.0}
\newcommand{\staticSame}{11}
\newcommand{\staticTotal}{12}
\newcommand{\staleJumpMin}{0.66}
\newcommand{\staleJumpMax}{1.36}
\newcommand{\staleViolAtBzero}{yes}
\newcommand{\foldMaxGap}{0.35}
\newcommand{\foldBfourWorst}{yes}
\newcommand{\realSpread}{0.014}
\newcommand{\snapAbiGap}{0.14}
\newcommand{\snapAbiChg}{0.19}
\newcommand{\bfourAbiGap}{0.014}
\newcommand{\boneGapPP}{0.007}
\newcommand{\maxAFfour}{1.117}
\newcommand{\maxAFlira}{1.0024}
\newcommand{\maxAFothers}{1.000}
\newcommand{\driftNobSwNone}{0.41}
\newcommand{\driftNobSwHot}{0.82}
\newcommand{\driftMaxCostPct}{0.74}
\newcommand{\dpRunAbi}{3.6}
\newcommand{\daggerRoundZero}{56}
\newcommand{\daggerRuns}{60}
\newcommand{\nobCand}{51}
\newcommand{\nobhlOchg}{0.80}
\newcommand{\nobJBtwobase}{875.3}
\newcommand{\nobJLiRAbase}{875.9}
\newcommand{\nobJBsixbase}{875.3}
\newcommand{\nobJObase}{874.7}
\newcommand{\nobJBtwohl}{882.2}
\newcommand{\nobJLiRAhl}{880.9}
\newcommand{\nobJBsixhl}{887.0}
\newcommand{\nobJOhl}{880.9}
\newcommand{\nobJBtwohd}{883.2}
\newcommand{\nobJLiRAhd}{883.7}
\newcommand{\nobJBsixhd}{883.7}
\newcommand{\nobJOhd}{882.9}
\newcommand{\nobJBtwost}{895.4}
\newcommand{\nobJLiRAst}{898.6}
\newcommand{\nobJBsixst}{899.4}
\newcommand{\nobJOst}{894.3}
\newcommand{\nobstGainPct}{0.09}
\newcommand{\liraHold}{2\,159}
\newcommand{\nobTestIntervals}{2\,160}
\newcommand{\hlLiraOracleDiff}{0}
\newcommand{\wearCredit}{0.040}
\newcommand{\sleepSaveabi}{0.060}
\newcommand{\sleepSavegea}{0.025}
\newcommand{\sleepSavenob}{0.035}
\newcommand{\effBsixbase}{0.07}
\newcommand{\effBsixDaysbase}{38}
\newcommand{\effBsixhl}{-0.75}
\newcommand{\effBsixDayshl}{50}
\newcommand{\nwJBfourbase}{1582.5}
\newcommand{\nwJLiRAbase}{1566.1}
\newcommand{\nwJStaticbase}{1566.1}
\newcommand{\nwJBtwobase}{1566.0}
\newcommand{\nwJBsixbase}{1566.0}
\newcommand{\nwJObase}{1565.9}
\newcommand{\nwBfourExcessbase}{1.1}
\newcommand{\nwJBfourhl}{1572.4}
\newcommand{\nwJLiRAhl}{1570.2}
\newcommand{\nwJStatichl}{1569.6}
\newcommand{\nwJBtwohl}{1570.0}
\newcommand{\nwJBsixhl}{1571.1}
\newcommand{\nwJOhl}{1569.6}
\newcommand{\nwBfourExcesshl}{0.2}
\newcommand{\nwJBfourhd}{1590.3}
\newcommand{\nwJLiRAhd}{1573.9}
\newcommand{\nwJStatichd}{1573.9}
\newcommand{\nwJBtwohd}{1574.0}
\newcommand{\nwJBsixhd}{1573.9}
\newcommand{\nwJOhd}{1573.9}
\newcommand{\nwBfourExcesshd}{1.0}
\newcommand{\nwJBfourst}{1584.8}
\newcommand{\nwJLiRAst}{1583.7}
\newcommand{\nwJStaticst}{1584.1}
\newcommand{\nwJBtwost}{1583.1}
\newcommand{\nwJBsixst}{1583.5}
\newcommand{\nwJOst}{1583.0}
\newcommand{\nwBfourExcessst}{0.1}
\newcommand{\nwSettings}{12}
\newcommand{\nwOconst}{9}
\newcommand{\nwLiraBest}{0}
\newcommand{\nwStaticSame}{10}
\newcommand{\nwHoldMaxGap}{0.12}
\newcommand{\nzarABfourChg}{57.8}
\newcommand{\nzarAOChg}{0.00}
\newcommand{\nzarALiRAChg}{0.00}
\newcommand{\nzarAExcess}{32.5}
\newcommand{\nzarALiRAgap}{0.00}
\newcommand{\nzarABsixgap}{0.00}
\newcommand{\nzarAStaticgap}{0.00}
\newcommand{\nzarABtwogap}{0.00}
\newcommand{\nzarBBfourChg}{1.9}
\newcommand{\nzarBOChg}{0.93}
\newcommand{\nzarBLiRAChg}{1.47}
\newcommand{\nzarBExcess}{0.6}
\newcommand{\nzarBLiRAgap}{0.36}
\newcommand{\nzarBBsixgap}{0.35}
\newcommand{\nzarBStaticgap}{0.36}
\newcommand{\nzarBBtwogap}{0.31}
\newcommand{\nzhalfBfourChg}{0.0}
\newcommand{\nzhalfOChg}{0.00}
\newcommand{\nzhalfLiRAChg}{0.00}
\newcommand{\nzhalfExcess}{0.0}
\newcommand{\nzhalfLiRAgap}{0.00}
\newcommand{\nzhalfBsixgap}{0.00}
\newcommand{\nzhalfStaticgap}{0.00}
\newcommand{\nzhalfBtwogap}{0.01}
\newcommand{\nzMainExcess}{6.5}
\newcommand{\nzMinP}{0.03}


\definecolor{cBlue}{RGB}{52,101,164}
\definecolor{cRed}{RGB}{196,78,82}
\definecolor{cGreen}{RGB}{85,153,96}
\definecolor{cOrange}{RGB}{221,132,52}
\definecolor{cPurple}{RGB}{129,114,179}
\definecolor{cGray}{RGB}{120,120,120}
\tikzset{box/.style={draw=#1!70!black,fill=#1!12,rounded corners=3pt,align=center,font=\footnotesize,minimum height=1.05cm,text width=2.85cm,line width=0.6pt},box/.default=cBlue,arr/.style={-{Stealth[length=2.2mm]},line width=0.6pt,draw=black!75},darr/.style={arr,dashed}}
\pgfplotsset{every axis/.append style={font=\footnotesize,grid=major,grid style={gray!25},legend style={font=\scriptsize,draw=gray!50,fill=white,fill opacity=0.9,text opacity=1}}}
\algrenewcommand\algorithmicrequire{\textbf{Input:}}
\algrenewcommand\algorithmicensure{\textbf{Output:}}

\renewcommand{\thesection}{S\arabic{section}}
\renewcommand{\thetable}{S\arabic{table}}
\renewcommand{\thefigure}{S\arabic{figure}}
\renewcommand{\thealgorithm}{S\arabic{algorithm}}
\begin{document}
\begin{center}{\Large Supplementary Material}\\[6pt]
When Does Reconfiguration Awareness Pay Off? An Exact-Oracle Study of Learned Link Sleeping in Software-Defined Networks\\[6pt]
Adeel Ahmad, Toqeer Ali Syed, Ali Akarma\end{center}

This file holds the algorithms, the secondary experiments, and the tables and figures that complement the article. Equation, section, and table numbers without the prefix S refer to the article. All values are generated from the experiment files in \url{https://github.com/adeliusa486/-lira-sdn-energy}.

\section{Related Work}\label{sec:related-full}

\subsection*{Energy-Aware Networking and Sleep Modes}

Surveys of green networking classify techniques into adaptive link rate, sleeping, and energy-aware routing~\cite{bianzino2012survey,bolla2011energy}. Chiaraviglio \emph{et al.} formulated ISP energy minimization as an integer program and proposed heuristics that switch off nodes and links subject to capacity and connectivity constraints~\cite{chiaraviglio2012minimizing}. GreenTE maximizes the number of links that can sleep under utilization and delay constraints~\cite{zhang2010greente}, ElasticTree applies the same principle to data-center networks~\cite{heller2010elastictree}, and EATe spreads traffic online so that lightly loaded links can sleep~\cite{vasic2010eate}. Power models for such studies rely on measured device profiles~\cite{mahadevan2009power,restrepo2009energy,vishwanath2014modeling}. In SDN, energy-aware routing has been formulated as integer programs with heuristics for carrier-grade Ethernet~\cite{maaloul2018carrier}, and a survey of SDN-specific methods is given in~\cite{assefa2019survey}.

\subsection*{Learning for Energy-Efficient Networking}

MER-SDN uses machine learning to predict the parameters of an energy-aware routing heuristic~\cite{assefa2018mersdn}. Reinforcement learning has been applied in an SDN controller to trade packet delay against energy~\cite{frohlich2021rl} and to compute energy-saving routes in data-center networks~\cite{zheng2022energysaving}. In cellular networks, deep reinforcement learning combined with traffic forecasting controls base-station sleep modes~\cite{wu2021deepbsc}, a problem that shares the on-off structure of link sleeping. Deep reinforcement learning has also been used to put switches and links of data-center networks to sleep~\cite{wang2022dcndrl} and to route traffic in SDN~\cite{pei2024drlte}, and energy-aware traffic engineering has been adapted to hybrid SDN deployments~\cite{maity2022ethos}. A recent study of two European ISP networks shows that more than a third of the links could sleep without congestion~\cite{rollin2024sleep}, which confirms that the opportunity measured on research backbones exists in operational networks. The closest works to ours replace the optimizer with a learned mapping from traffic matrices to precomputed configurations and report large speedups~\cite{gomezdelahiz2025noms,gomezdelahiz2026tgcn}. We reimplement this snapshot approach as baseline B4 with our own library, features, and classifier, because we did not have access to the original code. The comparison therefore concerns the snapshot principle and not the exact published systems.

\subsection*{Learning to Optimize and Learned Traffic Engineering}

Machine learning for combinatorial optimization treats problem instances as data and learns to imitate expert decisions or to replace parts of a solver~\cite{bengio2021mlco}. In traffic engineering, learning-accelerated and decomposition-based solvers~\cite{xu2023teal,perry2023dote,abuzaid2021ncflow,narayanan2021pop}, learned routing~\cite{valadarsky2017learning}, deep reinforcement learning with graph neural networks~\cite{almasan2022enero,bernardez2023magnneto}, and neural models of network performance~\cite{rusek2020routenet} have shown that learned components can approach solver quality at much lower latency. Most of these methods do not price the change from one configuration to the next. FlexDATE is an exception: it limits the disturbance of rerouting in reinforcement-learning traffic engineering~\cite{ye2023flexdate}, but it does not put links to sleep. LiRA instead imitates an oracle that is exact for the sequential objective, which the finite library makes tractable offline.

\subsection*{Reconfiguration Cost, Device Lifetime, and Switching Costs}

Consistent-update abstractions~\cite{reitblatt2012abstractions} and update scheduling~\cite{jin2014dionysus} show that changing forwarding state is neither free nor instantaneous. Bonetto \emph{et al.} included reconfiguration costs in multi-period power-aware logical topology design~\cite{bonetto2013multiperiod}. Chiaraviglio \emph{et al.} modeled how sleep modes affect device lifetime through temperature and power-state transitions~\cite{chiaraviglio2015lifetime,chiaraviglio2014impact}, and lifetime-aware provisioning has been studied for optical backbones~\cite{wiatr2016lifetime}. Online problems with switching costs are studied as metrical task systems~\cite{borodin1992mts} and smoothed online convex optimization~\cite{chen2018obd}, with dynamic right-sizing of data centers as a well-known application~\cite{lin2013rightsizing}. Our setting is a metrical task system with a finite state space whose offline optimum is computable, which makes imitation learning natural~\cite{ross2011dagger}. Learning-augmented online algorithms combine a learned predictor with a robust online rule, with guarantees for caching~\cite{lykouris2021caching}, general metrical task systems~\cite{antoniadis2023mts,christianson2023dart}, and smoothed online optimization~\cite{rutten2023smoothed}. They show that following a predictor blindly can be costly, which is the flapping we measure, and that a rule which holds its state is the natural robust fallback. Our hysteresis and static baselines are simple instances of such rules. Concept-drift adaptation~\cite{gama2014drift,lu2019drift} with detectors such as ADWIN~\cite{bifet2007adwin} and cumulative-sum tests~\cite{page1954cusum} provides the tools for deciding when to retrain, and drift is known to degrade learned models in network security~\cite{andresini2021insomnia}.



\section{Algorithms}\label{sec:algorithms}

Algorithm~\ref{alg:offline} summarizes offline training and Algorithm~\ref{alg:online} the online loop. In the emulation (Section~\ref{sec:res-emu}), reconfigurations follow a make-before-break order: links that the new configuration needs are woken first, new rules are installed, and only then are unused links put to sleep. The flow-level simulation charges each change through $R$ and $W$ and does not model the transient.





\begin{algorithm}[h!]
\caption{LiRA Offline Training}
\label{alg:offline}
\footnotesize
\begin{algorithmic}[1]
\Require training traces $D_{1:T}$, validation traces, prices $\lambda_R,\lambda_W$, rounds $N$
\Ensure library $\mathcal{C}$, policy $\pi$
\State $\mathcal{C}\gets$ \Call{BuildLibrary}{$D_{1:T}$} \Comment{ILP, NSGA-II, set cover}
\State precompute $E(c,D_t)$, feasibility, and $R$, $W$ for all pairs
\State compute $Q_t(c',c)$ and $V_t(c)$ by recursion (8)
\State $\mathcal{D}\gets\{(s_t,c^\star_t)\}$ on the oracle path ($s_t$ contains $c^\star_{t-1}$), weights $w$
\State train $\pi$ on $\mathcal{D}$
\For{$n=1$ to $N$}
\State roll out $\pi$ with the safety layer and expert mixing $\beta_n$; record $(t,c_{t-1})$
\State relabel the visited pairs from $Q_t(c_{t-1},\cdot)$ and add them to $\mathcal{D}$
\State retrain $\pi$ on $\mathcal{D}$
\EndFor
\State \Return $\mathcal{C}$ and the policy with the lowest validation cost
\end{algorithmic}
\end{algorithm}


\begin{algorithm}[h!]
\caption{LiRA Online Control Loop}
\label{alg:online}
\footnotesize
\begin{algorithmic}[1]
\Require $\mathcal{C}$, $\pi$, threshold $\theta$, initial $c_0$
\For{each interval $t=1,2,\dots$}
\State build $s_t$ from the traffic estimate, $c_{t-1}$, and its dwell time
\State $c\gets\pi(s_t)$
\If{$\mathrm{MLU}(c,\hat D_t)>\theta$}
\State $c\gets$ fallback rule (9)
\EndIf
\If{$c\neq c_{t-1}$}
\State wake links in $x\setminus x_{t-1}$; install new rules
\State sleep links in $x_{t-1}\setminus x$
\EndIf
\State $c_t\gets c$; observe $D_t$; update detectors with $\rho_t$ from \eqref{eq:proxy}
\If{a detector fires}
\State schedule retraining on the last seven days (Section~\ref{sec:drift})
\EndIf
\EndFor
\end{algorithmic}
\end{algorithm}


\section{Limitations and Threats to Validity}\label{sec:limitations-full}

\begin{itemize}
\item \emph{Traffic.} The only series with strong flapping is semi-synthetic. We vary its noise structure (Section~\ref{sec:res-noise}), but no real trace with such fluctuations was available to us. The evaluation on real traffic covers two research and education backbones, Abilene (2004) and G\'{E}ANT (2005), and both traces are old. The Nobel-Germany test split covers only \nobTestDays\ days, so tests on per-day costs have little power.
\item \emph{Power and wear.} The port power values and the wear parameters are assumed and not measured. The conclusion that the optimal configuration is stable depends on idle port power dominating dynamic power. A router with stronger load-proportional power, such as the high-dynamic-power scenario scaled further, may change this.
\item \emph{Baselines.} B4 is our reimplementation of the snapshot principle, not the published systems. B2 is restricted to the library. We did not solve the full multi-period mixed-integer program, but the per-interval bound limits what it can gain over the oracle (Section~6.3). B1, B3, and B5 were evaluated on subsamples. We did not train a deep reinforcement learning baseline.
\item \emph{Information.} The main results let every online method see the traffic of the interval it configures. With the previous interval's traffic, costs are unchanged but the safety layer admits violations on Nobel-Germany (Section~\ref{sec:res-stale}). We did not evaluate traffic forecasts or estimation error from link loads.
\item \emph{Drift.} The injected shifts did not change the optimal configuration, and the detectors fired independently of them, so the drift study is inconclusive. Link failures and large library refreshes were not evaluated.
\item \emph{Emulation.} The emulation is small (17 switches, one day, compressed time), uses a modeled wake-up delay, single-path routing per demand, and no real line cards, and it reports a single run for LiRA.
\item \emph{Statistics.} Seeds vary only the learner. The traffic traces are fixed and the per-day tests use few days, so no comparison survives a correction over all settings.
\item \emph{Learned component.} In \liraConstRuns\ of \liraTotalRuns\ LiRA runs the trained policy is constant, in \daggerRoundZero\ of them the selected policy predates DAgger, and the static rule equals LiRA in \staticSame\ of \staticTotal\ settings. A learned policy that does more than hold a configuration is therefore evaluated only in the Nobel-Germany stress scenario (\nobJLiRAst\ against \nobJOst\ for the oracle, five seeds). On G\'{E}ANT the classifier is not constant as a model, but its output is constant on the test split.
\end{itemize}



\section{Drift Detection and Retraining}\label{sec:drift}

The oracle cannot be computed online, but a myopic regret proxy can. After each decision, the controller computes
\begin{equation}\label{eq:proxy}
\begin{split}\rho_t=&E(c_t,D_t)-\min_{c\in\mathcal{C}:\,\mathrm{MLU}(c,D_t)\le\theta}E(c,D_t)\\
&+\mathbb{1}[\mathrm{MLU}(c_t,D_t)>\theta],\end{split}
\end{equation}
the energy lost against the best feasible configuration plus a violation indicator. ADWIN~\cite{bifet2007adwin} monitors $\rho_t$, and a Page--Hinkley test~\cite{page1954cusum} monitors the deseasonalized log of the total traffic volume, that is, the log volume minus its training mean for the same time of week. When a detector fires, the policy is retrained on the last seven days with oracle labels computed on that window. Because the window just after a change still contains mostly old data, retraining is repeated one and three days later. A library refresh, which adds integer-program solutions for eight clusters of the recent window, is available as an option.



\section{Dependence on the Snapshot Classifier}\label{sec:res-models}

The amount of flapping depends strongly on the classifier. On Nobel-Germany in the base scenario, B4 with XGBoost changes configuration \modXgbChg\ times per day, LightGBM \modLgbmChg, a random forest \modRfChg, and a multilayer perceptron \modMlpChg\ (Table~\ref{tab:models}). The random forest is nearly as stable as LiRA. Averaging over trees is a plausible reason, which we did not test. The perceptron flaps so much that its cost (\modMlpJ) is far above that of LiRA (\nobJlira). LiRA reaches the same cost with every backend, because its policy is constant in this scenario. A sequence-aware method is therefore most useful when the snapshot classifier is the kind that flaps, and it is least useful when a smoother classifier or a hysteresis guard is available.




\begin{table}[h!]
\centering
\caption{Choice of Model on Nobel-Germany (Base Scenario, Seed 0). Chg/d: Changes per Day}
\label{tab:models}
\footnotesize
\setlength{\tabcolsep}{3pt}
\begin{adjustbox}{max width=\linewidth}\begin{tabular}{@{}lccccc@{}}
\toprule
 & \multicolumn{3}{c}{B4 (snapshot)} & \multicolumn{2}{c}{LiRA} \\
\cmidrule(lr){2-4}\cmidrule(lr){5-6}
Model & Chg/d & $J$ & Latency (ms) & Chg/d & $J$ \\
\midrule
XGBoost & 12.80 & 937.2 & 1.09 & 0.27 & 875.9 \\
LightGBM & 10.13 & 924.1 & 1.61 & 0.27 & 875.9 \\
Random forest & 0.53 & 877.2 & 17.22 & 0.27 & 875.9 \\
MLP ($2\times256$) & 44.80 & 1093.7 & 0.72 & 0.27 & 875.9 \\
\bottomrule
\end{tabular}
\end{adjustbox}
\end{table}


\section{Ablation, Library Size, and Threshold}\label{sec:res-ablation}

Table~\ref{tab:ablation} removes one component of LiRA at a time on Nobel-Germany at the default price. It uses seeds 0 to 2, so the full policy's stress cost differs from the five-seed value in Table~2. In the high-load scenario the expert is nearly constant, so all variants that keep the safety layer reach the oracle cost. In the stress scenario the full policy costs \ablStressFull\ against \ablStressO\ for the oracle. Removing the previous-configuration input raises the cost to \ablStressNoPrev, removing the cost-sensitive weights to \ablStressNoCS, and removing DAgger to \ablStressNoDagger. Removing the lag features lowers it to \ablStressNoLag. The differences are small and of mixed sign, and we did not test them for significance. Removing the safety layer lowers the cost to \ablStressNoSafe\ because the policy then ignores the bound, and it causes \ablStressNoSafeViol\% of the intervals to violate $\theta$. Because $J$ carries no penalty for violations, this cost is not comparable, and the layer is needed. On Abilene and G\'{E}ANT all variants coincide with the oracle because the expert is constant.


Table~\ref{tab:K} shows the effect of the library size on Nobel-Germany. Requested sizes above \KactualMax\ yield the same library, because the candidate pool contains only \nobCand\ configurations. Every library size covers all test matrices. A larger library lowers the cost of B4 (from \KBFouratEight\ at $K=8$ to \KBFouratSixFour\ at the largest size), but hysteresis and LiRA stay within 0.2\% of the oracle cost at every size. Table~\ref{tab:theta} varies the utilization threshold. At $\theta=0.6$ B4 changes configuration \thSixBFourchg\ times per day and LiRA \thSixLiRAchg\ times. For $\theta\ge0.7$, each method gives the same result at every threshold on Nobel-Germany, because the same configurations are optimal and feasible. Table~\ref{tab:theta} uses seed 0, so its B4 change rate differs from the five-seed mean of Table~1. On Abilene and G\'{E}ANT, the threshold has no effect between 0.6 and 0.9. LiRA never violates the bound when the library covers the traffic.





\begin{table}[h!]
\centering
\caption{Ablation of LiRA on Nobel-Germany at the Default Price (Mean over Seeds 0 to 2)}
\label{tab:ablation}
\scriptsize
\setlength{\tabcolsep}{2pt}
\begin{adjustbox}{max width=\linewidth}\begin{tabular}{@{}lcccc@{}}
\toprule
Variant & $S$ (\%) & Changes/d & MLU viol. (\%) & $J$ \\
\midrule
\multicolumn{5}{@{}l}{\textit{Nobel-Germany, high load}} \\
Oracle O$^\star$ & 27.57 & 0.80 & 0.00 & 880.9 \\
LiRA (full) & 27.57 & 0.80 & 0.00 & 880.9 \\
no lag features & 27.57 & 0.80 & 0.00 & 880.9 \\
no previous config. & 27.57 & 0.80 & 0.00 & 880.9 \\
no cost-sensitive weights & 27.57 & 0.80 & 0.00 & 880.9 \\
no DAgger & 27.57 & 0.80 & 0.00 & 880.9 \\
no safety layer & 27.58 & 0.00 & 0.14 & 877.3 \\
\midrule
\multicolumn{5}{@{}l}{\textit{Nobel-Germany, stress}} \\
Oracle O$^\star$ & 27.35 & 0.80 & 0.00 & 894.3 \\
LiRA (full) & 27.37$\pm$0.03 & 0.93$\pm$0.33 & 0.00 & 897.4$\pm$2.2 \\
no lag features & 27.36$\pm$0.03 & 0.84$\pm$0.19 & 0.00 & 895.6$\pm$5.0 \\
no previous config. & 27.40 & 1.07 & 0.00 & 902.0 \\
no cost-sensitive weights & 27.37$\pm$0.03 & 1.07 & 0.00 & 898.6$\pm$7.8 \\
no DAgger & 27.37$\pm$0.03 & 0.89$\pm$0.38 & 0.00 & 899.0$\pm$9.1 \\
no safety layer & 27.40$\pm$0.07 & 0.00 & 0.17$\pm$0.07 & 889.9$\pm$1.6 \\

\bottomrule
\end{tabular}
\end{adjustbox}
\end{table}


\begin{table}[h!]
\centering
\caption{Library Size on Nobel-Germany: Test Coverage and Total Cost $J$ of Each Method (Base Scenario, Seed 0)}
\label{tab:K}
\footnotesize
\setlength{\tabcolsep}{4pt}
\begin{adjustbox}{max width=\linewidth}\begin{tabular}{@{}cccccc@{}}
\toprule
$K$ & Test coverage (\%) & B4 & B6 & LiRA & O$^\star$ \\
\midrule
8 & 100.0 & 939.3 & 875.3 & 875.9 & 874.7 \\
16 & 100.0 & 938.6 & 875.3 & 875.9 & 874.7 \\
32 & 100.0 & 937.2 & 875.3 & 875.9 & 874.7 \\
64 (51) & 100.0 & 935.7 & 875.3 & 875.8 & 874.7 \\
128 (51) & 100.0 & 935.7 & 875.3 & 875.8 & 874.7 \\
\bottomrule
\end{tabular}
\end{adjustbox}
\end{table}


\begin{table*}[h!]
\centering
\caption{Utilization Threshold $\theta$ (Base Scenario, Seed 0)}
\label{tab:theta}
\footnotesize
\setlength{\tabcolsep}{4pt}
\begin{adjustbox}{max width=\linewidth}\begin{tabular}{@{}lccccccc@{}}
\toprule
Topology & $\theta$ & All-on viol. (\%) & $S_{\mathrm{O}^\star}$ (\%) & B4 chg/d & LiRA chg/d & LiRA viol. (\%) & LiRA $J$ \\
\midrule
Abilene & 0.6 & 0.00 & 23.98 & 0.00 & 0.00 & 0.00 & 7226.7 \\
Abilene & 0.7 & 0.00 & 23.98 & 0.00 & 0.00 & 0.00 & 7226.7 \\
Abilene & 0.8 & 0.00 & 23.98 & 0.00 & 0.00 & 0.00 & 7226.7 \\
Abilene & 0.9 & 0.00 & 23.98 & 0.00 & 0.00 & 0.00 & 7226.7 \\
\midrule
G\'{E}ANT & 0.6 & 0.00 & 29.85 & 0.00 & 0.00 & 0.00 & 603.7 \\
G\'{E}ANT & 0.7 & 0.00 & 29.85 & 0.00 & 0.00 & 0.00 & 603.7 \\
G\'{E}ANT & 0.8 & 0.00 & 29.85 & 0.00 & 0.00 & 0.00 & 603.7 \\
G\'{E}ANT & 0.9 & 0.00 & 29.85 & 0.00 & 0.00 & 0.00 & 603.7 \\
\midrule
Nobel-Germany & 0.6 & 0.09 & 27.63 & 10.40 & 0.40 & 0.00 & 876.0 \\
Nobel-Germany & 0.7 & 0.05 & 27.63 & 12.80 & 0.27 & 0.00 & 875.9 \\
Nobel-Germany & 0.8 & 0.05 & 27.63 & 12.80 & 0.27 & 0.00 & 875.9 \\
Nobel-Germany & 0.9 & 0.05 & 27.63 & 12.80 & 0.27 & 0.00 & 875.9 \\

\bottomrule
\end{tabular}
\end{adjustbox}
\end{table*}


\section{Wear and Lifetime}\label{sec:res-wear}

Fig.~\ref{fig:wear} shows the mean acceleration factor of the decision sequences on Nobel-Germany for $\alpha=0.5$ and varying cycling weight $\chi$ (the sequences are fixed, and only the model changes). Every method stays below the always-on value of one for $\alpha=0.5$ and all $\chi$ up to $0.1$, but B4 degrades fastest: at $\chi=0.1$ its mean factor is \wearBigChiBFour, against \wearBigChiLiRA\ for LiRA. At the default parameters (seed 0) the most-cycled link of B4 has a factor of \maxAFfour, whereas the most-cycled link of LiRA has \maxAFlira\ and those of B6, B7, and the oracle \maxAFothers. When sleeping brings no benefit ($\alpha=1$), every transition lowers lifetime, and every method that switches exceeds the always-on factor, B4 by the largest margin. Absolute lifetime figures depend on assumed parameters and should be read as relative comparisons.




\begin{figure}[h!]
\centering
\begin{tikzpicture}
\begin{axis}[width=\columnwidth,height=4.2cm,xmode=log,xlabel={Cycling weight $\chi$ (per transition per day)},ylabel={Mean $\mathrm{AF}$},legend pos=north west,legend columns=2,ymin=0.8,ymax=1.05]
\addplot[cRed,thick,mark=*,mark size=1.5pt] table[x=chi,y=B4,col sep=space]{generated/wear_chi.dat};\addlegendentry{B4}
\addplot[cGreen!70!black,thick,mark=square*,mark size=1.5pt] table[x=chi,y=B6,col sep=space]{generated/wear_chi.dat};\addlegendentry{B6}
\addplot[cBlue,thick,dashed,mark=triangle*,mark size=1.5pt] table[x=chi,y=LiRA,col sep=space]{generated/wear_chi.dat};\addlegendentry{LiRA}
\addplot[black,dotted,thick] table[x=chi,y=B0,col sep=space]{generated/wear_chi.dat};\addlegendentry{Always on}
\end{axis}
\end{tikzpicture}
\caption{Mean acceleration factor of the decision sequences on Nobel-Germany for $\alpha=0.5$. The factor of the always-on network is one.}
\label{fig:wear}
\end{figure}


\section{Decision Time and Scalability}\label{sec:res-time}

The library-based decisions are fast. The median decision time of B4 is \abiLatBFour\,ms on Abilene, \nobLatBFour\,ms on Nobel-Germany, and \geaLatBFour\,ms on G\'{E}ANT. LiRA needs less than 0.01\,ms where its policy is constant and \geaLatLiRA\,ms where a tree model is evaluated. The safety check takes at most \safetyTA\,ms even on a 65-node topology. Solving the integer program per traffic matrix (B1) takes a median of \nobIlpMs\,ms on Nobel-Germany and \geaIlpS\,s on G\'{E}ANT, NSGA-II (B3) takes \nobGaS\,s and \geaGaS\,s, and the greedy heuristic is faster but gives away part of the savings (Table~\ref{tab:opt}). On the subsample, the per-snapshot integer program reaches the same savings as the library snapshot optimum to within \boneGapPP\ percentage points. On larger SNDlib topologies the integer program takes \ilpGermany\,s for Germany50 and \ilpTA\,s for TA2 (Table~\ref{tab:scale}). These times fit five-minute intervals on the topologies we study, so speed alone does not require a learned controller here. The dynamic-programming oracle needs \dpRunAbi\,s for all \abiIntervals\ Abilene intervals with $K=32$, so prices can be explored offline.





\begin{table}[h!]
\centering
\caption{Per-Snapshot Optimization Baselines on a Test Subsample}
\label{tab:opt}
\footnotesize
\setlength{\tabcolsep}{4pt}
\begin{adjustbox}{max width=\linewidth}\begin{tabular}{@{}llccc@{}}
\toprule
Topology & Method & $n$ & $S$ (\%) & Median time (s) \\
\midrule
Abilene & Library snapshot optimum & 60 & 23.98 & -- \\
Abilene & B1 per-snapshot ILP & 60 & 23.98 & 0.53 \\
Abilene & B3 NSGA-II & 15 & 23.98 & 253.97 \\
Abilene & B5 greedy & 60 & 23.55 & 0.04 \\
\midrule
G\'{E}ANT & Library snapshot optimum & 60 & 29.85 & -- \\
G\'{E}ANT & B1 per-snapshot ILP & 60 & 29.85 & 3.53 \\
G\'{E}ANT & B3 NSGA-II & 60 & 29.56 & 41.56 \\
G\'{E}ANT & B5 greedy & 60 & 27.16 & 0.24 \\
\midrule
Nobel-Germany & Library snapshot optimum & 60 & 27.63 & -- \\
Nobel-Germany & B1 per-snapshot ILP & 60 & 27.63 & 0.54 \\
Nobel-Germany & B3 NSGA-II & 60 & 27.63 & 28.34 \\
Nobel-Germany & B5 greedy & 60 & 27.29 & 0.09 \\

\bottomrule
\end{tabular}
\end{adjustbox}
\end{table}


\begin{table}[h!]
\centering
\caption{Scalability of the Optimization Step and the Safety Check}
\label{tab:scale}
\footnotesize
\setlength{\tabcolsep}{2pt}
\begin{adjustbox}{max width=\linewidth}\begin{tabular}{@{}lcccccc@{}}
\toprule
Topology & $|V|$ & $|E|$ & OD pairs & ILP (s) & Greedy (s) & Safety check (ms) \\
\midrule
Abilene & 12 & 15 & 132 & 0.48 & 0.04 & 0.067 \\
Nobel-Germany & 17 & 26 & 272 & 0.41 & 0.19 & 0.130 \\
G\'{E}ANT & 22 & 36 & 462 & 4.07 & 0.56 & 0.157 \\
Germany50 & 50 & 88 & 2450 & 1.61 & 1.45 & 0.223 \\
TA2 & 65 & 108 & 4160 & 11.37 & 7.42 & 0.363 \\
\bottomrule
\end{tabular}
\end{adjustbox}
\end{table}


\section{Without the Wear Price}\label{sec:res-nowear}

With the default prices, a sleeping link earns a wear credit about as large as its energy saving (Section~5), so $J$ is well below the energy of the network. To check that this assumed credit does not drive the comparison, we repeat the main study with $\lambda_W=0$: oracle labels, tuning of B6 and B7, training of B4 and LiRA, and the reported cost then contain energy and reconfiguration only. Table~\ref{tab:nowear} shows the result for all \nwSettings\ settings. The oracle holds a single configuration in \nwOconst\ of them. On Nobel-Germany B4 is again the most expensive online method (\nwBfourExcessbase\% above the oracle in the base scenario), while B6, B2, the static rule, and LiRA lie within \nwHoldMaxGap\% of the oracle in every setting. LiRA is cheaper than every other online method in \nwLiraBest\ settings and costs the same as the static rule in \nwStaticSame. Under high load it is slightly more expensive than the static rule (\nwJLiRAhl\ against \nwJStatichl, which equals the oracle). The ranking of the methods therefore does not depend on the wear model.




\begin{table*}[h!]
\centering
\caption{Repetition of the Main Study Without a Wear Price ($\lambda_W=0$, All Methods Retrained): Configuration Changes per Day and Cost $J$ (Energy plus Reconfiguration). B4, B6, and LiRA: Mean over Five Seeds}
\label{tab:nowear}
\footnotesize
\setlength{\tabcolsep}{3pt}
\begin{adjustbox}{max width=\linewidth}\begin{tabular}{@{}llcccccccccc@{}}
\toprule
 & & \multicolumn{4}{c}{Changes/d} & \multicolumn{6}{c}{Cost $J$ with $\lambda_W=0$} \\
\cmidrule(lr){3-6}\cmidrule(lr){7-12}
Topology & Scenario & B4 & B6 & LiRA & O$^\star$ & B4 & B6 & B2 & Static & LiRA & O$^\star$ \\
\midrule
Abilene & base & 0.02 & 0.00 & 0.00 & 0.00 & 9151.3 & 9151.0 & 9154.7 & 9151.0 & 9151.0 & 9151.0 \\
Abilene & high load & 0.00 & 0.00 & 0.00 & 0.00 & 9157.3 & 9157.3 & 9161.0 & 9157.3 & 9157.3 & 9157.3 \\
Abilene & high dynamic power & 0.02 & 0.00 & 0.00 & 0.00 & 9172.3 & 9172.0 & 9183.1 & 9172.0 & 9172.0 & 9172.0 \\
Abilene & stress & 0.00 & 0.00 & 0.00 & 0.00 & 9190.9 & 9190.9 & 9201.5 & 9190.9 & 9190.9 & 9190.9 \\
G\'{E}ANT & base & 0.00 & 0.00 & 0.00 & 0.00 & 1897.8 & 1897.8 & 1897.8 & 1897.8 & 1897.8 & 1897.8 \\
G\'{E}ANT & high load & 0.00 & 0.00 & 0.00 & 0.00 & 1904.4 & 1904.4 & 1904.4 & 1904.4 & 1904.4 & 1904.4 \\
G\'{E}ANT & high dynamic power & 0.00 & 0.00 & 0.00 & 0.00 & 1919.1 & 1919.1 & 1919.1 & 1919.1 & 1919.1 & 1919.1 \\
G\'{E}ANT & stress & 0.00 & 0.00 & 0.00 & 0.00 & 1938.9 & 1938.9 & 1938.9 & 1938.9 & 1938.9 & 1938.9 \\
Nobel-Germany & base & 11.68 & 0.13 & 0.27 & 0.00 & 1582.5 & 1566.0 & 1566.0 & 1566.1 & 1566.1 & 1565.9 \\
Nobel-Germany & high load & 1.76 & 1.07 & 0.88 & 0.80 & 1572.4 & 1571.1 & 1570.0 & 1569.6 & 1570.2 & 1569.6 \\
Nobel-Germany & high dynamic power & 11.68 & 0.27 & 0.27 & 0.27 & 1590.3 & 1573.9 & 1574.0 & 1573.9 & 1573.9 & 1573.9 \\
Nobel-Germany & stress & 1.76 & 1.07 & 0.93 & 0.80 & 1584.8 & 1583.5 & 1583.1 & 1584.1 & 1583.7 & 1583.0 \\
\bottomrule
\end{tabular}
\end{adjustbox}
\end{table*}


\section{Traffic Drift}\label{sec:res-drift}

We inject three shifts into the test period of each topology: a ramp that raises all demands by up to 50\% over ten days, a fourfold increase of three large demands, and a six-hour shift of the daily pattern. Table~\ref{tab:drift} compares a LiRA policy trained once with drift-triggered retraining. On Abilene and G\'{E}ANT the shifts did not change which configuration this policy uses, because the library stays feasible. On Nobel-Germany the hotspot raised its configuration changes from \driftNobSwNone\ to \driftNobSwHot\ per day. The cost changed by at most \driftMaxCostPct\% relative to the unshifted stream. The detectors were also not informative: without any shift, they raised \driftFaAbi, \driftFaGea, and \driftFaNob\ false alarms per 1000 intervals on Abilene, G\'{E}ANT, and Nobel-Germany, and the Page--Hinkley threshold calibrated on validation data was at the maximum of the search grid (\driftThrAbi) on the two real traces. We therefore cannot attribute detections to the injected shifts, and retraining did not lower the cost of LiRA. The injected shifts did not change the optimal configuration, so this experiment does not validate the detectors.




\begin{table}[h!]
\centering
\caption{Injected Traffic Shifts: Post-Shift Cost $J$ of LiRA Trained Once (``Once'') and with Drift-Triggered Retraining (``Retrained''). FA/1k Is the Number of False Alarms per 1000 Intervals on the Stream Without a Shift}
\label{tab:drift}
\footnotesize
\setlength{\tabcolsep}{3pt}
\begin{adjustbox}{max width=\linewidth}\begin{tabular}{@{}llcccc@{}}
\toprule
Topology & Shift & $J$ once & $J$ retrained & Detections & FA/1k \\
\midrule
Abilene & growth & 4701.5 & 4701.5 & 2 & 0.17 \\
Abilene & hotspot & 4698.9 & 4698.9 & 2 & 0.17 \\
Abilene & phase & 4698.2 & 4698.2 & 2 & 0.17 \\
\midrule
G\'{E}ANT & growth & 394.8 & 394.8 & 1 & 0.00 \\
G\'{E}ANT & hotspot & 392.6 & 392.6 & 0 & 0.00 \\
G\'{E}ANT & phase & 392.3 & 392.3 & 0 & 0.00 \\
\midrule
Nobel-Germany & growth & 569.9 & 570.6 & 2 & 0.93 \\
Nobel-Germany & hotspot & 573.9 & 574.5 & 2 & 0.93 \\
Nobel-Germany & phase & 569.6 & 570.3 & 2 & 0.93 \\

\bottomrule
\end{tabular}
\end{adjustbox}
\end{table}


\section{Rolling-Origin Folds}\label{sec:res-folds}

Table~\ref{tab:folds} repeats the comparison for four rolling origins, retraining B4, B6, and LiRA on all data before each origin. In every fold B4 costs more than B6 and LiRA, and B6 and LiRA stay within \foldMaxGap\% of the oracle. On Abilene and G\'{E}ANT, all methods coincide in every fold. On Nobel-Germany, the mean cost over the folds is \foldNobBFour\ for B4, \foldNobBSix\ for B6, \foldNobBSeven\ for B7, \foldNobLiRA\ for LiRA, and \foldNobO\ for the oracle.




\begin{table}[h!]
\centering
\caption{Rolling-Origin Folds (Mean over Four Folds, Default Price)}
\label{tab:folds}
\footnotesize
\setlength{\tabcolsep}{3pt}
\begin{adjustbox}{max width=\linewidth}\begin{tabular}{@{}llccc@{}}
\toprule
Topology & Method & $S$ (\%) & Chg/d & $J$ (mean$\pm$std over folds) \\
\midrule
Abilene & B4 & 23.98 & 0.00 & 2890.5$\pm$0.9 \\
 & B6 & 23.98 & 0.00 & 2890.5$\pm$0.9 \\
 & LiRA & 23.98 & 0.00 & 2890.5$\pm$0.9 \\
 & O$^\star$ & 23.98 & 0.00 & 2890.5$\pm$0.9 \\
\midrule
G\'{E}ANT & B4 & 29.84 & 0.00 & 241.7$\pm$0.7 \\
 & B6 & 29.84 & 0.00 & 241.7$\pm$0.7 \\
 & LiRA & 29.84 & 0.00 & 241.7$\pm$0.7 \\
 & O$^\star$ & 29.84 & 0.00 & 241.7$\pm$0.7 \\
\midrule
Nobel-Germany & B4 & 27.64 & 11.17 & 371.5$\pm$5.4 \\
 & B6 & 27.64 & 0.25 & 350.1$\pm$0.5 \\
 & LiRA & 27.64 & 0.33 & 350.3$\pm$0.7 \\
 & O$^\star$ & 27.63 & 0.17 & 350.0$\pm$0.5 \\

\bottomrule
\end{tabular}
\end{adjustbox}
\end{table}


\section{Emulation of Configuration Changes}\label{sec:res-emu}

We emulate the most active day of the Nobel-Germany series in Mininet 2.3.0 with Open vSwitch 3.7.1 (17 switches, one host per switch, 26 links). Each configuration installs, for every origin--destination pair, the path that carries the largest share of its demand as OpenFlow rules with \texttt{ovs-ofctl}. A day is replayed with two seconds per interval. A configuration change wakes the links that the new configuration needs, waits one second of modeled wake-up time (Mininet does not model line-card physics), replaces the flow tables of the switches that change one after another, and then puts unused links to sleep. Eight UDP probe flows at 100 packets per second on the largest origin--destination pairs measure loss and one-way delay. B4 makes \emuTrans\ configuration changes on this day, each toggling two links, whereas the sequences of B6, LiRA, and the oracle make none. Each change of B4 takes \emuTransS\,s in total, which is the modeled wake-up time of 1\,s plus \emuRuleMs\,ms of measured rule installation and the time to put links to sleep. Table~\ref{tab:emu} reports the probe loss. Each run sends \emuProbePkts\ probe packets. B4 loses \emuLostPerTrans\ packets per configuration change on average, all of them inside the windows around a change: \emuLossWin\% of the probes in those windows and \emuLossAll\% over the whole day. The always-on control and LiRA lose no probe packets, and LiRA keeps the same link-hours as B4, computed from the replayed configurations, which are \emuLinkSave\% fewer than those of the always-on network. This emulation measures the cost of a single configuration change, not the energy of the network, and it ignores line-card boot time beyond the modeled delay.




\begin{table*}[h!]
\centering
\caption{Mininet Emulation of the Most Active Day of Nobel-Germany (B4: Three Runs; B0 and LiRA: One Run). Trans.: Configuration Changes. Loss in Windows Counts Probes Sent from 0.2\,s Before a Change Starts to 1\,s After It Ends. Link-Hours Are Computed from the Replayed Configurations}
\label{tab:emu}
\footnotesize
\setlength{\tabcolsep}{4pt}
\begin{adjustbox}{max width=\linewidth}\begin{tabular}{@{}lccccccccc@{}}
\toprule
Method & Runs & Trans. & Rule install (ms) & Trans. time (s) & Loss (\%) & Loss in windows (\%) & Delay p50 (ms) & p99 (ms) & Link-hours \\
\midrule
B0 & 1 & 0 & -- & -- & 0.000 & -- & 0.10 & 0.4 & 4.16 \\
B4 & 3 & 18 & 99 & 1.11 & 0.014 & 0.22 & 0.10 & 0.3 & 2.88 \\
LiRA & 1 & 0 & -- & -- & 0.000 & -- & 0.09 & 0.2 & 2.88 \\
\bottomrule
\end{tabular}
\end{adjustbox}
\end{table*}


\section{Cost of Ignoring the Sequence}

For two configurations $c_a$ and $c_b$, let $s_{ab}=\lambda_R R(c_a,c_b)+\lambda_W\chi\|x_a-x_b\|_1$ be the non-energy cost of switching between them.

\begin{proposition}\label{prop:snapshot}
Let $c_a,c_b\in\mathcal{C}$ be feasible in every interval and have the same number of sleeping links, and let the traffic alternate so that $E(c_a,D_t)-E(c_b,D_t)=(-1)^t\epsilon$ with $\epsilon>0$, while every other library member is dominated in every interval. Then the snapshot rule (7) incurs a regret of at least $(T-1)s_{ab}-\lceil T/2\rceil\epsilon$ with respect to the optimum of (6), which grows linearly in $T$ whenever $s_{ab}>\epsilon/2$.
\end{proposition}
\begin{proof}
The snapshot rule selects $c_b$ in even and $c_a$ in odd intervals, so it pays the lower energy in each interval and switches $T-1$ times. The sleep terms of (5) are identical for $c_a$ and $c_b$, so the two sequences differ only in energy and in switching cost. Staying at $c_a$ throughout pays at most $\lceil T/2\rceil\epsilon$ more energy and never switches. Its cost is an upper bound on the optimum, which proves the claim.
\end{proof}

The proposition shows that flapping can be costly in principle. Whether it is costly in practice depends on whether real traffic produces such near-ties, which is the empirical question of Section~6.



\section{Implementation and Parameters}\label{sec:params}

The flow-level simulator is written in Python with NumPy and sparse path--link matrices. Integer and linear programs are solved with HiGHS (highspy 1.15.1). Experiments ran on an Intel Core i9-13900H laptop (14 cores, 20 threads, 48\,GB RAM, no GPU) under Windows, and the emulation ran in Ubuntu on WSL2 with Mininet 2.3.0 and Open vSwitch 3.7.1. Table~\ref{tab:params} lists the parameters. The port power values are representative assumptions within the ranges reported in~\cite{vishwanath2014modeling,mahadevan2009power}, not measurements of a specific device, and the wear parameters are assumed values for the model of~\cite{chiaraviglio2015lifetime}. We therefore study the sensitivity to both (Section~\ref{sec:res-wear}).




\begin{table}[h!]
\centering
\caption{Model and Training Parameters}
\label{tab:params}
\footnotesize
\setlength{\tabcolsep}{3pt}
\begin{tabular}{@{}p{3.9cm}p{4.6cm}@{}}
\toprule
Parameter & Value \\
\midrule
Threshold $\theta$; paths $k$; library $K$ & 0.8; 4; 32 (default) \\
Port power $P^{\mathrm{idle}}$/$P^{\mathrm{dyn}}$/$P^{\mathrm{sl}}$ & 30/5/3\,W \\
Wear $\alpha$; $\chi$ & 0.5; 0.01 per transition per day \\
Prices $\lambda_R$; $\lambda_W$ & $0.002\,n_d$; $0.0008\,n_d/\chi$ \\
Interval $\Delta$; $n_d$ & 5\,min, 288 (G\'{E}ANT: 15\,min, 96) \\
XGBoost & depth 6, rate 0.1, 150 rounds max, early stop 15 \\
Features & PCA 95\% variance, lags 1/3/6 \\
DAgger rounds; $\beta_1$; weight clip & 3; 0.5; $[0.2,5]$ \\
ADWIN confidence; Page--Hinkley tolerance, forgetting & 0.002; 0.005, 0.9999 \\
Retraining window; schedule & 7 days; at detection, $+1$ and $+3$ days \\
Seeds; folds & 5; 4 rolling origins \\
\bottomrule
\end{tabular}
\end{table}


\section{Noise Structure of the Semi-Synthetic Series}\label{sec:res-noise}

The only series with flapping is semi-synthetic, and its noise is independent across intervals. To see how the findings depend on this choice, we generate three further 30-day series with the same gravity model, weekly profile, split, and scaling. Two use log-normal noise with the same marginal scale but AR(1) correlation in time ($\rho=0.5$ and $\rho=0.9$), and one uses independent noise with half the scale. Each variant gets its own library, and every method is retrained (base scenario, five seeds). Table~\ref{tab:noise} shows that the amount of flapping depends strongly on the noise. With $\rho=0.5$ the snapshot classifier changes configuration \nzarABfourChg\ times per day and costs \nzarAExcess\% more than the oracle, because demands now stay near a decision boundary for several intervals and cross it repeatedly. With $\rho=0.9$ the oracle itself switches (\nzarBOChg\ changes per day) and B4 is only \nzarBExcess\% above it, and with half the noise B4 does not flap at all. In every variant, B6, B2, the static rule, and LiRA are at most \nzarBLiRAgap\% above the oracle, and LiRA costs the same as the static rule. The value of reconfiguration awareness thus ranges from nothing to a third of the cost, depending on the traffic, and in none of the variants does it require learning.




\begin{table*}[h!]
\centering
\caption{Nobel-Germany Series with Different Noise (Base Scenario): Configuration Changes per Day, Total Cost $J$, and the Excess Cost of B4 over the Oracle. B4, B6, and LiRA: Mean over Five Seeds}
\label{tab:noise}
\footnotesize
\setlength{\tabcolsep}{3pt}
\begin{adjustbox}{max width=\linewidth}\begin{tabular}{@{}lccccccccccc@{}}
\toprule
 & \multicolumn{4}{c}{Changes/d} & \multicolumn{6}{c}{Total cost $J$} & B4 above \\
\cmidrule(lr){2-5}\cmidrule(lr){6-11}
Noise & B4 & B6 & LiRA & O$^\star$ & B4 & B6 & B2 & Static & LiRA & O$^\star$ & O$^\star$ (\%) \\
\midrule
i.i.d., $\sigma=1.23$ (main) & 11.68 & 0.13 & 0.27 & 0.00 & 931.7 & 875.3 & 875.3 & 875.9 & 875.9 & 874.7 & 6.5 \\
AR(1) $\rho=0.5$, $\sigma=1.23$ & 57.81 & 0.00 & 0.00 & 0.00 & 1158.5 & 874.4 & 874.4 & 874.4 & 874.4 & 874.4 & 32.5 \\
AR(1) $\rho=0.9$, $\sigma=1.23$ & 1.92 & 1.47 & 1.47 & 0.93 & 890.4 & 888.0 & 887.7 & 888.1 & 888.1 & 884.9 & 0.6 \\
i.i.d., $\sigma=0.62$ & 0.00 & 0.00 & 0.00 & 0.00 & 884.2 & 884.2 & 884.3 & 884.2 & 884.2 & 884.2 & 0.0 \\
\bottomrule
\end{tabular}
\end{adjustbox}
\end{table*}


\section{Decisions on the Previous Interval's Traffic}\label{sec:res-stale}

All results so far let the controller see $D_t$ when it configures interval $t$ (Section~5.3). Table~\ref{tab:info} repeats the Nobel-Germany comparison when every online method decides with $D_{t-1}$: the features, the safety check, the hysteresis and dwell decisions, and the forecast of B2 all use the last completed interval, while costs and violations are measured on $D_t$ (the same trained models, seeds 0 to 2). The decisions barely change, and costs move by at most \staleNobMaxDJ\% on Nobel-Germany and by \staleRealDJtext\% on Abilene and G\'{E}ANT. The utilization guarantee does not survive. On Nobel-Germany, B2, B4, and B6 violate $\theta$ in \staleBzerobase\%, \staleBzerohl\%, and \staleBzerost\% of the intervals in the base, high-load, and stress scenarios, the same shares as the always-on network with shortest-path routing, and LiRA in \staleLiRAViolbase\%, \staleLiRAViolhl\%, and \staleLiRAViolst\%. Every violated interval of the static rule is one in which the always-on network also violates $\theta$, and in each of them the utilization of the installed configuration rises by \staleJumpMin\ to \staleJumpMax\ within one interval. Such peaks come from the independent per-interval noise of the series and cannot be predicted from $D_{t-1}$. A margin of $0.05$ on $\theta$ therefore does not remove them, and it raises the cost (LiRA under high load: \marginLiRAhlJ\ against \nobJLiRAhl). On Abilene and G\'{E}ANT at most \staleRealMaxViol\% of the intervals violate $\theta$. LiRA's exact match with the oracle under high load thus depends on knowing the traffic of the interval it configures.




\begin{table}[h!]
\centering
\caption{Nobel-Germany: Total Cost $J$ and Share of Intervals Violating $\theta$ When Every Online Method Knows $D_t$ or Decides with $D_{t-1}$, Without and With a Safety Margin of 0.05 (B4, B6, and LiRA: Mean over Seeds 0 to 2)}
\label{tab:info}
\scriptsize
\setlength{\tabcolsep}{3pt}
\begin{adjustbox}{max width=\linewidth}\begin{tabular}{@{}lccccc@{}}
\toprule
 & $D_t$ known & \multicolumn{2}{c}{$D_{t-1}$} & \multicolumn{2}{c}{$D_{t-1}$, margin 0.05} \\
\cmidrule(lr){2-2}\cmidrule(lr){3-4}\cmidrule(lr){5-6}
Method & $J$ & $J$ & Viol. (\%) & $J$ & Viol. (\%) \\
\midrule
\multicolumn{6}{@{}l}{\textit{Nobel-Germany, base}} \\
B2 rolling horizon & 875.3 & 875.3 & 0.05 & -- & -- \\
B4 snapshot ML & 933.7 & 933.7 & 0.05 & 933.7 & 0.05 \\
B6 hysteresis & 875.3 & 875.3 & 0.05 & 875.3 & 0.05 \\
Static + fallback & 875.9 & 875.8 & 0.05 & 875.8 & 0.05 \\
LiRA & 875.9 & 875.8 & 0.05 & 875.8 & 0.05 \\
Oracle O$^\star$ & 874.7 & -- & -- & -- & -- \\
\midrule
\multicolumn{6}{@{}l}{\textit{Nobel-Germany, high load}} \\
B2 rolling horizon & 882.2 & 882.2 & 0.19 & -- & -- \\
B4 snapshot ML & 891.2 & 891.2 & 0.19 & 897.1 & 0.19 \\
B6 hysteresis & 887.0 & 887.0 & 0.19 & 890.7 & 0.19 \\
Static + fallback & 880.9 & 880.9 & 0.14 & 886.2 & 0.14 \\
LiRA & 880.9 & 880.9 & 0.14 & 886.2 & 0.14 \\
Oracle O$^\star$ & 880.9 & -- & -- & -- & -- \\
\midrule
\multicolumn{6}{@{}l}{\textit{Nobel-Germany, stress}} \\
B2 rolling horizon & 895.4 & 895.3 & 0.19 & -- & -- \\
B4 snapshot ML & 903.7 & 903.7 & 0.19 & 909.6 & 0.19 \\
B6 hysteresis & 899.4 & 899.4 & 0.19 & 903.9 & 0.19 \\
Static + fallback & 902.0 & 902.0 & 0.19 & 909.6 & 0.19 \\
LiRA & 897.4 & 897.4 & 0.17 & 901.4 & 0.17 \\
Oracle O$^\star$ & 894.3 & -- & -- & -- & -- \\
\bottomrule
\end{tabular}
\end{adjustbox}
\end{table}


\section{Notation, Parameters, and Additional Figures}
\begin{table}[h!]
\centering
\caption{Main Notation}
\label{tab:notation}
\footnotesize
\setlength{\tabcolsep}{3pt}
\begin{tabular}{@{}p{1.85cm}p{6.5cm}@{}}
\toprule
Symbol & Meaning \\
\midrule
$G=(V,E)$ & Network graph; $E$ is the set of undirected links \\
$\kappa_e$ & Capacity of link $e$ in each direction \\
$\Delta$, $n_d$ & Control interval; number of intervals per day \\
$D_t$, $\hat D_t$ & Traffic matrix of interval $t$; its estimate \\
$c=(x,r)$ & Configuration: link states $x\in\{0,1\}^{|E|}$, routing $r$ \\
$\mathcal{C}$, $K$ & Configuration library and its size \\
$u_e$, $\theta$ & Utilization of link $e$; utilization threshold \\
$E(c,D)$ & Energy of $c$ under traffic $D$ in one interval \\
$R(c',c)$ & Reconfiguration cost of moving from $c'$ to $c$ \\
$\mathrm{AF}_e$, $\alpha$, $\chi$ & Acceleration factor; sleep factor; cycling weight \\
$W(c',c)$ & Per-step wear cost \\
$\lambda_R,\lambda_W$ & Prices of reconfiguration and wear \\
$J$, $V_t$, $Q_t$ & Total cost; cost-to-go; action value \\
\bottomrule
\end{tabular}
\end{table}

\begin{figure}[h!]
\centering
\begin{tikzpicture}
\begin{axis}[width=\columnwidth,height=4.2cm,xlabel={Day of the busiest test week},ylabel={Configuration index},ymin=-0.3,ymax=3.4,legend pos=north east,legend columns=3,mark size=0.8pt]
\addplot[const plot,cRed,thick] table[x=day,y expr=\thisrow{B4}+0.0,col sep=space]{generated/case_week.dat};
\addlegendentry{B4}
\addplot[const plot,cGreen!70!black,thick] table[x=day,y expr=\thisrow{B6}+0.12,col sep=space]{generated/case_week.dat};
\addlegendentry{B6 (+0.12)}
\addplot[const plot,cBlue,thick,dashed] table[x=day,y expr=\thisrow{LiRA}+0.24,col sep=space]{generated/case_week.dat};
\addlegendentry{LiRA (+0.24)}
\end{axis}
\end{tikzpicture}
\caption{Installed configuration index over the busiest test week of the Nobel-Germany series. Offsets separate the lines. B4 switches repeatedly between two configurations, whereas B6 and LiRA hold a configuration.}
\label{fig:case}
\end{figure}

\begin{figure}[h!]
\centering
\begin{tikzpicture}
\begin{groupplot}[group style={group size=2 by 1,horizontal sep=1.3cm},width=0.46\columnwidth,height=4.3cm,xmode=log,xlabel={Changes/day $+0.05$},scaled y ticks=false]
\nextgroupplot[title={Base},ylabel={Total cost $J$}]
\addplot[only marks,mark=square*,cGreen!70!black] table[x expr=\thisrow{config_changes_per_day}+0.05,y=total_cost_J,col sep=space]{generated/pareto_base_B6.dat};
\addplot[only marks,mark=triangle*,cOrange] table[x expr=\thisrow{config_changes_per_day}+0.05,y=total_cost_J,col sep=space]{generated/pareto_base_B7.dat};
\addplot[only marks,mark=*,cBlue,mark size=2pt] table[x expr=\thisrow{config_changes_per_day}+0.05,y=total_cost_J,col sep=space]{generated/pareto_base_LiRA.dat};
\addplot[only marks,mark=x,black,mark size=3pt] table[x expr=\thisrow{config_changes_per_day}+0.05,y=total_cost_J,col sep=space]{generated/pareto_base_Ostar.dat};
\addplot[only marks,mark=diamond*,cRed,mark size=3pt] table[x expr=\thisrow{config_changes_per_day}+0.05,y=total_cost_J,col sep=space]{generated/pareto_base_B4.dat};
\nextgroupplot[title={High load}]
\addplot[only marks,mark=square*,cGreen!70!black] table[x expr=\thisrow{config_changes_per_day}+0.05,y=total_cost_J,col sep=space]{generated/pareto_highload_B6.dat};
\addlegendentry{B6 sweep}
\addplot[only marks,mark=triangle*,cOrange] table[x expr=\thisrow{config_changes_per_day}+0.05,y=total_cost_J,col sep=space]{generated/pareto_highload_B7.dat};
\addlegendentry{B7 sweep}
\addplot[only marks,mark=*,cBlue,mark size=2pt] table[x expr=\thisrow{config_changes_per_day}+0.05,y=total_cost_J,col sep=space]{generated/pareto_highload_LiRA.dat};
\addlegendentry{LiRA}
\addplot[only marks,mark=x,black,mark size=3pt] table[x expr=\thisrow{config_changes_per_day}+0.05,y=total_cost_J,col sep=space]{generated/pareto_highload_Ostar.dat};
\addlegendentry{O$^\star$}
\addplot[only marks,mark=diamond*,cRed,mark size=3pt] table[x expr=\thisrow{config_changes_per_day}+0.05,y=total_cost_J,col sep=space]{generated/pareto_highload_B4.dat};
\addlegendentry{B4}
\end{groupplot}
\end{tikzpicture}
\caption{Total cost versus configuration changes on Nobel-Germany. Hysteresis thresholds and dwell times are swept for B6 and B7 and the reconfiguration price is swept for LiRA and the oracle. Each marker is one setting on the test split. The horizontal axis is logarithmic, and 0.05 is added to the change rate so that methods without changes can be shown.}
\label{fig:pareto}
\end{figure}

\begin{figure}[h!]
\centering
\begin{tikzpicture}
\begin{axis}[width=\columnwidth,height=4.0cm,ybar,bar width=14pt,ymin=0,ymax=110,ylabel={Share of changes (\%)},xlabel={Energy saved by a change (\% of interval energy)},xtick=data,xticklabels={$<0.1$,0.1--0.5,0.5--1,1--2,$>2$},nodes near coords={\pgfmathprintnumber[fixed,precision=0]{\pgfplotspointmeta}},nodes near coords style={font=\scriptsize},enlarge x limits=0.15]
\addplot[fill=cRed!55,draw=cRed] table[x=bin,y=share,col sep=space]{generated/gain_hist.dat};
\end{axis}
\end{tikzpicture}
\caption{Energy gained by each configuration change of the snapshot classifier B4 (seed 0) on the Nobel-Germany test period (\gainChanges\ changes). \gainTenthShare\% of the changes save less than 0.1\% of the interval energy.}
\label{fig:gain}
\end{figure}

\begin{table*}[h!]
\centering
\caption{Topologies and Traffic Data}
\label{tab:data}
\footnotesize
\setlength{\tabcolsep}{2.5pt}
\begin{adjustbox}{max width=\linewidth}\begin{tabular}{@{}lcccccc@{}}
\toprule
Topology & $|V|$ & $|E|$ & $\Delta$ & Intervals & Span & Traffic \\
\midrule
Abilene & 12 & 15 & 5\,min & 48\,108 & 2004-03-01 to 2004-09-10 & real (SNDlib) \\
G\'{E}ANT & 22 & 36 & 15\,min & 10\,784 & 2005-05-04 to 2005-08-31 & real (SNDlib) \\
Nobel-Germany & 17 & 26 & 5\,min & 8\,640 & 2004-03-01 to 2004-03-30 & semi-synthetic \\
\bottomrule
\end{tabular}
\end{adjustbox}
\end{table*}

\bibliographystyle{elsarticle-num}
\bibliography{references}
\end{document}
